% Este capítulo puede editarse y compilarse de forma individual.
% Para compilar el libro completo, usa main.tex.
\documentclass[../main.tex]{subfiles}
\begin{document}
\chapter{Teorema de no-cloning}

Antes de estudiar el resultado formal, conviene separar dos ideas que en la mecánica clásica suelen confundirse. En teoría clásica, si se posee un sistema físico y se conoce su estado, no parece problemático imaginar una copia ideal de ese estado. En mecánica cuántica la situación cambia: un \textbf{estado puro desconocido} no es simplemente una lista de valores preexistentes que podamos leer y duplicar sin perturbar. El estado cuántico contiene amplitudes, fases relativas y relaciones de superposición que sólo se manifiestan estadísticamente al medir muchos sistemas preparados del mismo modo.

Este fue precisamente uno de los puntos que hizo fértil la discusión de Herbert y, poco después, de Wootters y Zurek~\cite{HerbertFLASH,WoottersZurek1982}. Si un estado cuántico arbitrario pudiera copiarse perfectamente, entonces se podría amplificar la información contenida en un solo sistema y medir muchas copias para reconstruir propiedades incompatibles. En el contexto de correlaciones tipo EPR, esa posibilidad amenazaría el principio de no señalización. El \textbf{teorema de no-cloning} muestra que tal \textbf{máquina universal no existe}.

\begin{figure}[htbp]
    \centering
    \begin{tikzpicture}[
    >=Latex,
    caja/.style={draw=UdeCAzul, fill=UdeCAzulClaro, rounded corners, semithick, align=center,
        text width=2.8cm, minimum height=1.15cm, inner sep=6pt, font=\small},
    bloque/.style={draw=UdeCAzul, fill=UdeCAzul!15, rounded corners, semithick, align=center,
        text width=2.3cm, minimum height=1.55cm, font=\small},
    flecha/.style={->, thick, UdeCAzul, shorten >=3pt, shorten <=3pt},
    lab/.style={font=\footnotesize, fill=white, inner sep=2pt}
]
    \node[caja] (psi) at (-4.7,1.0) {estado desconocido $\ket{\psi}_{\theta}$};
    \node[caja] (blank) at (-4.7,-1.0) {registro en blanco $\ket{A}_{c}$};
    \node[bloque] (U) at (0,0) {$U$\\¿copiador universal?};
    \node[caja] (out1) at (4.7,1.0) {registro original $\ket{\psi}_{\theta}$};
    \node[caja] (out2) at (4.7,-1.0) {registro copia $\ket{\psi}_{c}$};
    \draw[flecha] (psi.east) -- node[lab, above] {entrada} (U.160);
    \draw[flecha] (blank.east) -- (U.200);
    \draw[flecha] (U.20) -- node[lab, above] {salida} (out1.west);
    \draw[flecha] (U.340) -- (out2.west);
\end{tikzpicture}
    \caption{Esquema ideal de una máquina universal de clonado. La imposibilidad aparece al exigir que el mismo operador $U$ funcione para cualquier superposición desconocida.}
    \label{fig:maquina-no-cloning}
\end{figure}

Supongamos que se desea copiar un estado puro desconocido $\ket{\psi}_\theta$ usando un segundo sistema, inicialmente preparado en un estado fijo $\ket{A}_c$. Idealmente, el proceso tendría que actuar como
\begin{equation*}
    \ket{\psi}_\theta\ket{A}_c
    \longmapsto
    \ket{\psi}_\theta\ket{\psi}_c.
\end{equation*}
La palabra ``universal'' significa que el mismo procedimiento debe funcionar para todo estado $\ket{\psi}$; la palabra ``exacto'' significa que la copia final debe ser idéntica al original; y la palabra ``determinista'' significa que el procedimiento debe ocurrir con probabilidad uno. Para un sistema cerrado, esto se modela mediante una transformación unitaria $U$ tal que
\begin{equation*}
    U\left(\ket{\psi}_\theta\ket{A}_c\right)
    =
    \ket{\psi}_\theta\ket{\psi}_c,
    \qquad
    \forall\,\ket{\psi}_\theta.
\end{equation*}

Para fijar ideas tomemos dos qubits. El primero, denotado por $\theta$, contiene el estado que queremos copiar; el segundo, denotado por $c$, será el registro copia. Elegimos el estado inicial del registro copia como
\begin{equation*}
    \ket{A}_c=\ket{0}_c.
\end{equation*}
Si la máquina copia correctamente los estados de la base computacional, debe satisfacer
\begin{equation*}
    U\left(\ket{0}_\theta\ket{0}_c\right)
    =
    \ket{0}_\theta\ket{0}_c,
    \qquad
    U\left(\ket{1}_\theta\ket{0}_c\right)
    =
    \ket{1}_\theta\ket{1}_c.
\end{equation*}
Hasta aquí no hay contradicción: copiar dos estados ortogonales conocidos sí es compatible con la mecánica cuántica. El problema aparece cuando se exige copiar una superposición arbitraria
\begin{equation*}
    \ket{\psi}_\theta=\alpha\ket{0}_\theta+\beta\ket{1}_\theta,
    \qquad
    |\alpha|^2+|\beta|^2=1.
\end{equation*}
Por \textbf{linealidad} de $U$ se obtiene
\begin{align*}
    U\left[\left(\alpha\ket{0}_\theta+\beta\ket{1}_\theta\right)\ket{0}_c\right]
    &=U\left(\alpha\ket{0}_\theta\ket{0}_c+
    \beta\ket{1}_\theta\ket{0}_c\right)\\
    &=\alpha U\left(\ket{0}_\theta\ket{0}_c\right)
    +\beta U\left(\ket{1}_\theta\ket{0}_c\right)\\
    &=\alpha\ket{0}_\theta\ket{0}_c+
    \beta\ket{1}_\theta\ket{1}_c.
\end{align*}
Sin embargo, si el proceso hubiera clonado de verdad el estado $\ket{\psi}$, el resultado tendría que ser
\begin{align*}
    \ket{\psi}_\theta\ket{\psi}_c
    &=
    \left(\alpha\ket{0}_\theta+\beta\ket{1}_\theta\right)
    \left(\alpha\ket{0}_c+\beta\ket{1}_c\right)\\
    &=
    \alpha^2\ket{0}_\theta\ket{0}_c
    +\alpha\beta\ket{0}_\theta\ket{1}_c
    +\alpha\beta\ket{1}_\theta\ket{0}_c
    +\beta^2\ket{1}_\theta\ket{1}_c.
\end{align*}
\textbf{En general, estos dos resultados no coinciden}. Por tanto, el mismo operador que copia correctamente $\ket{0}$ y $\ket{1}$ no copia correctamente una superposición arbitraria de ellos.

\begin{teorema}[No-cloning]
No existe una transformación unitaria universal que, para todo estado puro desconocido $\ket{\psi}$, realice el mapa
\begin{equation*}
    \ket{\psi}\ket{A}\longmapsto \ket{\psi}\ket{\psi},
\end{equation*}
con $\ket{A}$ fijo e independiente de $\ket{\psi}$.
\end{teorema}

\begin{proof}
La demostración anterior muestra la contradicción usando explícitamente la linealidad. También existe una prueba compacta usando preservación de productos internos. Supongamos que la máquina copia dos estados normalizados arbitrarios $\ket{\psi}$ y $\ket{\varphi}$:
\begin{equation*}
    U\ket{\psi}\ket{A}=\ket{\psi}\ket{\psi},
    \qquad
    U\ket{\varphi}\ket{A}=\ket{\varphi}\ket{\varphi}.
\end{equation*}
Como $U$ es unitario, preserva productos internos. Luego
\begin{align*}
    \braket{\psi|\varphi}
    &=\braket{\psi|\varphi}\braket{A|A}\\
    &=\left(\bra{\psi}\bra{A}\right)U^\dagger U\left(\ket{\varphi}\ket{A}\right)\\
    &=\left(\bra{\psi}\bra{\psi}\right)\left(\ket{\varphi}\ket{\varphi}\right)\\
    &=\braket{\psi|\varphi}^2.
\end{align*}
Por tanto, si $s=\braket{\psi|\varphi}$, se tiene $s=s^2$. Esta condición sólo se satisface, en general, si $s=0$ o $s=1$. Es decir, la copia perfecta sólo puede exigirse simultáneamente para estados ortogonales o idénticos. No puede existir una máquina universal que copie todos los estados cuánticos puros.
\end{proof}

\begin{observacion}
El teorema no dice que sea imposible copiar información clásica codificada en estados ortogonales. Por ejemplo, si sólo se usan los estados $\ket{0}$ y $\ket{1}$, una compuerta tipo CNOT puede copiar el valor lógico en un registro auxiliar. Lo que está prohibido es copiar de manera exacta, determinista y universal un estado cuántico desconocido. En particular, no se puede convertir una sola partícula en muchas copias idénticas para reconstruir todas sus amplitudes y fases mediante tomografía.
\end{observacion}

\begin{comentario}
Existen máquinas de clonado aproximado y protocolos probabilísticos que funcionan bajo hipótesis más restrictivas. La palabra clave es que dejan de ser simultáneamente universales, exactos y deterministas. El resultado de Wootters y Zurek se refiere precisamente a la imposibilidad de la copia perfecta de un cuanto arbitrario.
\end{comentario}
\end{document}
